% PENDING integration/root audit. Written during checkpoint08 freeze.
\begin{saresult}{A.123}{Proved under stated hypotheses}{Trivial low-dimensional component groups and the ball case}
Let $P_{\rm PL}(B)$ fix $B\times\{0\}\cup\partial B\times I$
pointwise, with the upper face free. Then
\[
 \pi_0P_{\rm PL}(B)=0
\]
if $B$ is a compact manifold of dimension zero or one, if $B$ is
a PL ball of any dimension, if $B$ is a compact orientable
surface other than $S^2$, or if $B$ is a closed nonorientable
surface other than $\mathbb{RP}^2$. The same assertion holds for a finite
disjoint union of these manifolds and for their definable models
over every real closed field, in the stipulated component convention.
No higher homotopy-group assertion is made here. The two closed
surfaces $S^2$ and $\mathbb{RP}^2$, and nonorientable surfaces
with boundary, are not included in the surface clauses.
\end{saresult}
\begin{proof}
All isotopies in this proof are jointly PL. In particular,
replacing an isotopy parameter $u$ by a product $st$ is unnecessary.
We use instead the PL function
\[
 \rho(s,t)=\max\{0,s+t-1\},\qquad (s,t)\in I^2.
                                                               \tag{A.123.1}
\]
It is zero if $s=0$ or $t=0$, equals $t$ if $s=1$, and equals
$s$ if $t=1$. Composition with a finite PL family therefore
remains a finite PL map, while preserving every parameter that
is explicitly retained as an output coordinate.

First suppose $B=D^m$. A concordance $h$ restricts on the upper
face to a PL homeomorphism $a$ of $D^m$ fixed on its boundary.
The PL Alexander isotopy gives a jointly PL path $a_u$ from the
identity to $a^{-1}$, fixed on $\partial D^m$; the explicit
level-preserving cone construction is given in
\cite[Theorem~5.2, pp.~95--96]{Epstein1966}.
For completeness, for $m\geq1$ on the cube $D^m=[-1,1]^m$ it is
\[
 a_u(x)=
 \begin{cases}
 u\,a^{-1}(x/u),&\|x\|_\infty\leq u,\ u>0,\\
 x,&\|x\|_\infty\geq u,
 \end{cases}
 \qquad a_0(x)=x.
\]
The two formulas agree on their common boundary. If
$a^{-1}(y)=Ay+b$ on a cell, the first formula is $Ax+ub$
on the cone over that cell at height one. These are finitely
many polyhedral cells, so the total isotopy is PL in every
dimension; the same formula using $a$ gives its inverse.
The maps
\[
 A_s(x,t)=(a_{\rho(s,t)}(x),t)
\]
are concordances, start at the identity, and make $A_1h$ the
identity on the entire boundary of $D^m\times I$.
This cylinder is a PL ball, so a second relative Alexander
isotopy deforms $A_1h$ to the identity. Concatenating the two
isotopies proves the ball assertion, including $m=0$ and $m=1$.
The displayed map and its inverse are continuous and PL: each
slice is a homeomorphism, and the total map is a continuous
bijection of compact polyhedra. Equivalently the inverse is
obtained by substituting the inverse PL isotopy family.

It remains in dimension one to treat $B=S^1=\mathbb R/\mathbb Z$.
Write $A=S^1\times I$. A concordance is orientation preserving
and its upper-circle restriction has degree one, since it is
homotopic to the bottom restriction through the cylinder.
An orientation-preserving PL circle homeomorphism is joined to
the identity by a PL isotopy: lift it to an increasing PL map
$\widetilde a:\mathbb R\to\mathbb R$ satisfying
$\widetilde a(x+1)=\widetilde a(x)+1$, triangulate the finitely
many breakpoint rectangles in $[0,1]\times I$, and linearly
interpolate their bottom and top vertex values. The horizontal
slopes on the two triangles of each rectangle are respectively
the positive slopes of $\widetilde a$ and of the identity.
Thus every horizontal slice is increasing, and the periodic
boundary values agree. This gives the required circle isotopy.
Using (A.123.1) extends an inverse upper-circle isotopy across
$A$, fixed on its bottom. We may consequently assume $h$ fixes
both boundary circles.

Let $\alpha=\{0\}\times I$. The proper arc $h(\alpha)$ has
an integral winding number $k$ relative to its fixed endpoints,
computed in the strip cover $\mathbb R\times I$.
The PL homeomorphism
\[
 T_{-k}([x],t)=([x-kt],t)
\]
fixes both boundary circles and makes $T_{-k}h(\alpha)$
homotopic to $\alpha$ relative to the endpoints. Moreover
$T_{-k}$ lies in the identity component of the concordance
group: its explicit isotopy is
\[
 T_{-k,s}([x],t)=([x-k\rho(s,t)],t).
\]
Only the bottom circle is required to stay fixed during this
isotopy. Epstein's relative arc theorem
\cite[Theorem~3.1, p.~89]{Epstein1966} supplies a PL ambient
isotopy fixed on both boundary circles taking the adjusted arc
to $\alpha$. Its interval parametrization can also be fixed,
using the relative interval Alexander isotopy in a regular
neighborhood, exactly as in the proof of
\cite[Theorem~3.4, p.~91]{Epstein1966}.
Cutting $A$ along $\alpha$ now gives a PL disk homeomorphism
fixed on its entire boundary. The relative Alexander isotopy
on this disk descends to the annulus. This proves the circle
assertion without retaining a Dehn-twist obstruction: the
upper face was free during the displayed $T_{-k,s}$ isotopy.

Finally let $F$ be a compact connected orientable surface
different from $S^2$. Waldhausen explicitly works in the PL
category, and in Section~3 assumes precisely that $F$ is
orientable and is not the sphere
\cite[pp.~56,65]{Waldhausen1968}. His Lemma~3.5, p.~69,
states that a homeomorphism of $F\times I$ fixed on
$F\times\{0\}\cup\partial F\times I$ is isotopic, keeping
the entire boundary fixed, to a level-preserving homeomorphism.
Apply the lemma to a concordance $h$, and write the resulting
homeomorphism as $(x,t)\mapsto(h_t(x),t)$, where $h_0$ is the
identity and every $h_t$ fixes $\partial F$.
Then
\[
 (x,t,s)\longmapsto(h_{\rho(s,t)}(x),t,s)
\]
is a jointly PL isotopy through concordances from the identity
to this level-preserving map. Its inverse is obtained from the
inverse level-preserving family. Composing with Waldhausen's
isotopy proves the surface clause, including the torus and all
orientable surfaces with boundary.

For completeness, closed nonorientable surfaces other than
$\mathbb{RP}^2$ also satisfy the assertion; this step uses the
nonorientable three-dimensional isotopy theorem, rather than
extending Waldhausen's orientable hypothesis without proof.
Let $F$ be such a surface and $V=F\times I$.
The universal cover of $F$ is the plane and its fundamental
group is torsion free (see the surface facts in
\cite[Corollary~1.8, pp.~85--86, and Lemma~4.3, p.~92]{Epstein1966}),
so the universal cover of $V$ is
$\mathbb R^2\times I$. An embedded PL two-sphere in $V$
lifts to a two-sphere bounding a PL ball in this slab, by the
three-dimensional PL Schoenflies theorem. Distinct translates
of its lifted sphere are disjoint. Their bounded balls are
either disjoint or nested. Proper discontinuity excludes
nesting: iteration of the corresponding nontrivial, infinite
order deck transformation would place infinitely many points
of an orbit in one compact ball. Thus the ball projects
injectively, and $V$ is irreducible. A two-sided embedded
$\mathbb{RP}^2$ cannot occur: its fundamental group maps
trivially to the torsion-free $\pi_1V$, so it would lift to
the orientable slab, where a two-sided embedded surface must
be orientable. Finally $F\times\{1/2\}$ is a two-sided
incompressible surface, since its inclusion is a homotopy
equivalence. It cannot bound a contractible submanifold:
otherwise its nontrivial injective fundamental-group map to
$V$ would factor through the zero fundamental group of that
submanifold. Hence $V$ is $\mathbb{RP}^2$-irreducible and
sufficiently large in exactly the definitions of
\cite[Chapter~II, Section~1, pp.~15--17]{Laudenbach1974}.

Write a concordance as $h(x,t)=(u(x,t),v(x,t))$.
It is properly homotopic to the identity. Indeed, first use
\[
 (x,t)\longmapsto
 \bigl(u(x,t),(1-s)v(x,t)+st\bigr),
\]
and then use
\[
 (x,t)\longmapsto\bigl(u(x,(1-s)t),t\bigr),\qquad 0\leq s\leq1.
\]
Here ``properly'' is the boundary-pair convention: throughout,
each of the two boundary components maps into itself, and
interior points remain interior. These formulas are only
homotopies of maps; they need not be isotopies.
Laudenbach's Theorem~II.6.1, pp.~46--47, deforms a proper
homotopy of a homeomorphism to the identity to an isotopy in
this class of manifolds. His category comparison in Chapter~I, Section~1,
pp.~7--9, first gives a TOP isotopy of $h$ to a smooth map $g$.
The map $g$ is still properly homotopic to the identity, so the
smooth formulation of Theorem~II.6.1 gives a TOP isotopy of
$h$ to the identity. The injectivity of
$\pi_0\operatorname{Aut}_{\rm PL}(V)\to
 \pi_0\operatorname{Aut}_{\rm TOP}(V)$ on pp.~7--8 then gives
a jointly PL isotopy $H_s:V\to V$ between the original PL
endpoints. This uses no assertion that the auxiliary smooth
map is literally PL in a fixed triangulation.

It remains to enforce the fixed bottom, since no unmentioned
relative version of that theorem is needed. Every $H_s$
preserves each of the two boundary components, by continuity
from its initial value. Let $b_s$ be its restriction to the
bottom, regarded as a PL homeomorphism of $F$.
The inverse total boundary family is PL, so
\[
 \widehat H_s=(b_s^{-1}\times\operatorname{id}_I)\,H_s
\]
is a jointly PL isotopy through concordances.
Both $b_0$ and $b_1$ are the identity. Thus this isotopy has
the same endpoints $h$ and the identity and fixes the bottom
pointwise at every time. This proves the additional closed
nonorientable clause, including the Klein bottle.

Fixing the bottom prevents permutation of components, so finite
disjoint unions give finite products of the groups just treated.
All real-field proofs consist of finite PL data. A.115 and the
relative component argument of A.113 transfer them to any real
closed field and identify these PL components with the
definable components. For manifolds with boundary the same
relative argument fixes the additional subpolyhedron
$\partial B\times I$ throughout.
\end{proof}
